\begin{chapintro}
同じ問いを調べた試験でも、結果は試験ごとにばらつきます。メタ解析は、そうした試験の結果を重みを付けて平均し、一つの試験よりも精度の高い答えを出す方法です。この章では、架空の降圧薬 A の 10 試験を例に、統合すると何が分かるのか、どこで誤るのかを説明し、この冊子全体で使う工程の全体図を示します。
\end{chapintro}

\section{一つの試験では答えが出ない理由}
\label{sec:01-1}

新しい降圧薬が脳卒中を減らすかどうかを知りたいとします。こうした問いには、ランダム化比較試験（RCT）で答えるのがふつうです。RCT は、参加者をくじ引きのような方法（ランダム化）で 2 つの群に分け、一方に調べたい薬剤を、もう一方にプラセボ（見た目は同じで有効成分を含まない錠剤）を投与して、その後の脳卒中の発症を比べる研究です。ランダム化によって、年齢や持病などの背景が 2 つの群で平均的にそろうので、発症の差を薬剤の作用によるものと考えることができます。

それでも、一つの試験の結果には偶然の揺れが含まれます。150 人ずつの小さな試験で、プラセボ群の脳卒中が 17 人、降圧薬群が 7 人だったとしても、同じ試験をもう一度行えば、17 人と 7 人ではなく 14 人と 10 人になるかもしれません。検査の内部精度管理で、同じ管理試料を測っても毎回わずかに違う値が出るのと同じで、標本が小さいほど揺れは大きくなります。

揺れの大きさを示すのが 95\% 信頼区間です。脳卒中の起こりやすさを群どうしで割り算した値をリスク比と呼び、リスク比が 1 なら差がなく、1 より小さければ降圧薬群で脳卒中が少ないことを表します。95\% 信頼区間は、同じ試験を何度も繰り返して毎回区間を計算したとき、そのおよそ 95\% が真の値を含むように作られた区間です。ただし、手元の 1 回の区間に真の値が 95\% の確率で入る、という意味ではありません。真の値は決まった一つの値で、区間のほうが試験のたびに動きます。95\% は、同じ手順を繰り返したときに区間が真の値を含む割合という、手順の性質を表しています。この区間が 1 をまたいでいれば、その試験だけでは「差がない」可能性を否定できません。

\section{降圧薬 A の 10 試験}
\label{sec:01-2}

表 1.1 は、この冊子の第 1 部を通して使う架空の 10 試験です。どの試験も高血圧の成人を対象に、降圧薬 A とプラセボを比べ、脳卒中の発症を数えています。対象が高齢者に限られた試験が 5 本、一般の成人を対象にした試験が 5 本あります。

\begin{kkcenter}
\kkfit{\begin{tabular}{lrrrrl}
\toprule
試験 & \multicolumn{2}{c}{降圧薬 A} & \multicolumn{2}{c}{プラセボ} & リスク比（95\% 信頼区間） \\
 & 事象 & 人数 & 事象 & 人数 & \\
\midrule
Aoki 2009 & 7 & 150 & 17 & 150 & 0.41（0.18〜0.96） \\
Berger 2011 & 48 & 1200 & 65 & 1200 & 0.74（0.51〜1.06） \\
Chen 2012 & 44 & 420 & 39 & 420 & 1.13（0.75〜1.70） \\
Dubois 2014 & 134 & 2500 & 145 & 2500 & 0.92（0.74〜1.16） \\
Endo 2015 & 6 & 90 & 12 & 90 & 0.50（0.20〜1.27） \\
Fischer 2016 & 32 & 800 & 38 & 800 & 0.84（0.53〜1.33） \\
Garcia 2018 & 121 & 1600 & 186 & 1600 & 0.65（0.52〜0.81） \\
Hayashi 2019 & 10 & 300 & 23 & 300 & 0.43（0.21〜0.90） \\
Ivanov 2021 & 112 & 3000 & 139 & 3000 & 0.81（0.63〜1.03） \\
Jensen 2022 & 43 & 650 & 83 & 650 & 0.52（0.36〜0.74） \\
\bottomrule
\end{tabular}}
\kkcap{表 1.1　架空の降圧薬 A の 10 試験（事象は脳卒中の発症）}
\end{kkcenter}

10 本のうち、95\% 信頼区間が 1 をまたがない試験、つまり単独で統計的に有意な差（偶然だけでは説明しにくい差）を示した試験は、Aoki 2009、Garcia 2018、Hayashi 2019、Jensen 2022 の 4 本だけです。残りの 6 本は区間が 1 をまたいでいて、Chen 2012 にいたっては点推定値が 1.13 と、降圧薬 A の群のほうが脳卒中が多くなっています。

この表だけを見て「有意な試験は 4 本、有意でない試験が 6 本なので、降圧薬 A が脳卒中を減らすかどうかは分からない」と結論するのは誤りです。この数え方は票数え（vote counting）と呼ばれ、メタ解析の教科書がそろって避けるよう勧めている方法です。有意でない試験の多くは、効果がないのではなく、人数が少なくて揺れが大きいために区間が広がり、1 をまたいでいるにすぎません。票数えでは、150 人ずつの試験と 3000 人ずつの試験が同じ 1 票になり、大きな試験のもつ情報が生かされません。

\section{統合すると何が分かるか}
\label{sec:01-3}

メタ解析では、各試験のリスク比を、試験の精度に応じた重みを付けて平均します。人数が多く事象も多い試験ほど推定が安定しているので、重みを大きくします。重みの付け方は \kkref{sec:05-3} で説明します。

10 試験を統合すると、リスク比は 0.73、95\% 信頼区間は 0.61〜0.86 になります。区間は 1 をまたがず、単独の試験よりもずっと狭くなりました。降圧薬 A の群では、脳卒中がプラセボ群のおよそ 4 分の 3 に減っていたと推定できます。プラセボ群で 100 人中 10 人が脳卒中を起こす集団なら、降圧薬 A によって 100 人中およそ 7 人に減る計算です。

10 試験の結果はかなりばらついています。リスク比は 0.41 から 1.13 まで散らばり、その散らばりは偶然の揺れだけで説明するには大きいと判定されました（検定の方法は \kkref{sec:06-1} で説明します）。高齢者の試験と一般の試験で効果が違うのか、試験の質の違いによるのか、それとも別の理由があるのかは、\kkref{sec:06-4} で調べます。統合した一つの数字を出して終わりにせず、ばらつきの理由を調べることも、メタ解析の仕事に含まれます。

\subsection{\hatten 累積メタ解析で結論の動きをたどる}
\label{sec:01-3-1}

試験を足していく順序を決めると、統合の結果が積み上がっていく様子が見えます。発表された年の順に 1 試験ずつ足し、足すたびに統合し直す方法を累積メタ解析と呼びます。図 1.1 の各行は、その行の試験までを足した時点の、統合したリスク比と 95\% 信頼区間です。

\kkfig[.9\linewidth]{x1_cumulative.pdf}{図 1.1　累積メタ解析。発表年の順に試験を足したときの、統合したリスク比と 95\% 信頼区間（架空のデータ）}

横軸は対数目盛のリスク比で、縦の実線が 1（差なし）です。対数目盛の意味は \kkref{sec:05-1-2} で説明します。一番上の行は Aoki 2009 だけの結果で、リスク比 0.41（0.18〜0.96）です。150 人ずつの小さな試験ですが、区間は 1 をまたいでいません。2 試験目の Berger 2011 を足すと 0.63（0.38〜1.05）、3 試験目の Chen 2012 を足すと 0.78（0.48〜1.25）と点推定が 1 に近づき、区間も 1 をまたぎます。大きな Dubois 2014 を足した 4 試験目では 0.87（0.71〜1.07）と区間が狭まりますが、まだ 1 をまたいでいます。区間は 6 試験目まで 1 をまたぎ続け、7 試験目の Garcia 2018 を足すと 0.78（0.65〜0.95）となって、以後は 1 をまたぎません。10 試験すべてでは 0.73（0.61〜0.86）です。

最初の小さな試験の 0.41 は、試験を足すうちに 0.73 まで上がりました。小さな試験が示す大きな効果は、偶然の揺れで大きく見えている可能性があります。また、有意な試験が 4 本しかなくても、大きな試験を足すたびに統合した区間が狭まり、7 試験目で 1 から離れました。ただし、この図は試験を発表年の順に足した 1 通りの結果にすぎません。累積メタ解析は、新しい試験が出るたびに結論がどう動いたかをたどる探索のための図で、区間が 1 をまたがなくなった時点を根拠に「何試験あれば十分だった」とは言えません。作り方は発展の \kkref{sec:07-7} で示します。

\section{メタ解析が答えられないこと}
\label{sec:01-4}

メタ解析は計算としては平均をとるだけなので、入れた試験が偏っていれば、出てくる答えも偏ります。この性質は「ごみを入れればごみが出る（garbage in, garbage out）」と言い表されます。偏りが生じる経路は大きく 3 つあります。

一つ目は、集めた試験そのものの質です。ランダム化の手順が不十分だった試験や、脱落者が多い試験は、それ自体が偏った結果を出しているおそれがあります。統合しても偏りは消えず、信頼区間が狭くなるぶん、偏った結果が確かなものに見えてしまいます。試験の質の評価は \kkref{sec:04-3} で扱います。

二つ目は、どの試験を集めたかです。効果が有意だった試験は論文になりやすく、有意でなかった試験は論文にならずに埋もれやすいことが知られています。発表された試験だけを集めると、効果を大きく見積もることになります。これを出版バイアスと呼び、調べ方は \kkref{sec:07-1} 以降で説明します。文献を漏れなく集める工夫（\kkref{sec:03-1}）は、この偏りを小さくするために行います。

三つ目は、違うものを混ぜることです。対象者も薬剤の量も追跡期間もまったく違う試験を一緒に平均すると、その平均値はどの集団にも当てはまらない数字になるおそれがあります。どこまでを同じ問いとみなして統合するかは、統計の計算を始める前に、問いの立て方（\kkref{sec:02-1}）で決めておきます。

検査の世界に置き換えると、測定法の違う施設のデータを、換算の確認もせずに一つの基準範囲にまとめるようなものです。数字は出ますが、その数字が何を意味するかは、元のデータの性質で決まります。

\section{系統的レビューの工程}
\label{sec:01-5}

メタ解析は、系統的レビューという作業の中の一工程です。系統的レビューは、あらかじめ書いた計画に沿って、関連する研究を漏れなく探し、決めた基準で選び、データを取り出し、質を評価したうえで結果をまとめる一連の作業を指します。統合の計算そのものは全体のごく一部で、作業時間の大半は文献の検索、選別、データの抽出に使われます。

図 1.2 は、この冊子で扱う工程を順に並べたものです。各工程の横に、その工程を説明する章の番号を付けました。第 2 部の検査のメタ解析も、統合の計算の部分を除けば、同じ工程をたどります。

\begin{kkcenter}
\begin{kkfitpic}
\begin{tikzpicture}[
  st/.style={kbox, minimum width=38mm, minimum height=7mm, font=\footnotesize},
  ch/.style={font=\scriptsize\color{kkdark}, inner sep=1mm},
  node distance=3.2mm]
\node[st] (a) {問いを立てる（PICO）};
\node[st, below=of a] (b) {計画を書き、登録する};
\node[st, below=of b] (c) {文献を検索する};
\node[st, below=of c] (d) {抄録と本文で選別する};
\node[st, below=of d] (e) {データを抽出する};
\node[st, right=12mm of a] (f) {バイアスのリスクを評価する};
\node[st, below=of f] (g) {効果を統合する};
\node[st, below=of g] (h) {ばらつきを調べる};
\node[st, below=of h] (i) {確かさを評価する（GRADE）};
\node[st, below=of i] (j) {報告する（PRISMA）};
\foreach \x/\y in {a/b,b/c,c/d,d/e,f/g,g/h,h/i,i/j} \draw[kar] (\x) -- (\y);
\draw[kar] (e.east) -- ++(6mm,0) |- (f.west);
\node[ch, anchor=east] at (a.west) {第 2 章};
\node[ch, anchor=east] at (b.west) {第 2 章};
\node[ch, anchor=east] at (c.west) {第 3 章};
\node[ch, anchor=east] at (d.west) {第 3 章};
\node[ch, anchor=east] at (e.west) {第 4 章};
\node[ch, anchor=west] at (f.east) {第 4 章};
\node[ch, anchor=west] at (g.east) {第 5 章};
\node[ch, anchor=west] at (h.east) {第 6 章};
\node[ch, anchor=west] at (i.east) {第 7 章};
\node[ch, anchor=west] at (j.east) {第 8 章};
\end{tikzpicture}
\end{kkfitpic}
\kkcap{図 1.2　系統的レビューとメタ解析の工程（左列の下から右列の上へ続く）}
\end{kkcenter}

工程は一方向に進むとは限りません。データを抽出している途中で、組み入れ基準の解釈が人によって違うことに気づき、選別に戻ることもあります。そうした変更は、いつ、なぜ行ったかを記録しておき、報告の際に書きます。後から基準を変えると結果が都合よく動いてしまうおそれがあるためで、計画を先に書いて登録しておく理由もここにあります（\kkref{sec:02-4}）。

\section{この章のまとめ}
\label{sec:01-6}

一つの試験の結果は偶然の揺れを含み、有意な試験の数を数える方法では答えを誤ります。メタ解析は試験の精度に応じた重みで平均をとる方法で、降圧薬 A の 10 試験からはリスク比 0.73（95\% 信頼区間 0.61〜0.86）という推定が得られました。ただし、統合の結果は集めた試験の質と範囲で決まるので、統合の前の工程を丁寧に行う必要があります。次の章では、その最初の工程として、問いを PICO の形に書き下し、計画を文書にするところから始めます。
